\cotefile{160}
\scanpage{13}
\[
 \Delta_n(\unsure{U|u}\mathbin{.}X)=(\det U)\,\Delta_n(X)\qquad\big]
\]
\inserted{$u\in\Gamma_n$}
\deleted{\unsure{autrement dit} \illegible{1} (\unsure{$U\in\Gamma_n$} \illegible{1} \unsure{ceci dit|coefficient})}

\deleted{donc les $\Delta_i=\Delta_n A_i$ sont aussi des invariants sous $S\Gamma_n$.}

\uline{Corollaire}\quad
\[
 \deleted{\mathbb F_p[X_1,\ldots,X_n]^{S\Gamma_n}\simeq
 \mathbb F_p[A_1,\ldots,\unsure{A_n|\Delta_n}]}
\]
\[
 \mathbb F_p[X_1,-,X_n]^{S\Gamma_n}
 =\mathbb F_p[A_1,\ldots,A_{n-1},\Delta_n]\quad
 \deleted{\unsure{\simeq}\ \mathbb F_p[X_1,\ldots,X_n}
\]
(et $\Gamma_n/S\Gamma_n\overset{\det}{\simeq} F_p^*$ opère sur
$\mathbb F_p[A_1,-,A_{n-1},\Delta_n]$ en laissant $A_1-A_{n-1}$ et par homothéties sur $\Delta_n$ ..

\uline{Dém} On a le diagramme commutatif
(\inserted{\unsure{flèches|flèche} $i^0$} défini par l'inclusion
$\mathbb F_p[A_1,-,A_{n-1},\unsure{\omega_n|W_n|\Delta_n}]\subset
\deleted{\mathbb F}\mathbb F_p[X_1,-,X_n]^{S\Gamma_n}$)
\diagram{\begin{tikzcd}[ampersand replacement=\&,row sep=large,column sep=large]
\mathbb E^n/\Gamma_n^{0}\arrow[r,"i^0"]\arrow[d,"p=\mathrm{can}" ] \& \mathbb E^n\arrow[d,"q"]\\
\mathbb E^n/\Gamma_n\arrow[r,"\sim"] \& \mathbb E^n
\end{tikzcd}}
$q$ définie par l'inclusion
$\mathbb F_p[A_1,-,A_{n-1},\pm\Delta_n^{p-1}]\subset
\mathbb F_p[A_1,-,A_{n-1},\Delta_n\unsure{]}$

Ainsi $p$ et $q$ sont des rev. galoisiens de groupe $\mathbb F_p^*$, identifiant le but au quotient de la source par ledit groupe, et $i^0$ est un \unsure{hom|homom.} compatible avec ces opérations, donc géom. bijectif et birationnel, donc (comme $\mathbb E^n$ est normal) un iso par le Main Theorem.
% archive sheet number: 12
